\documentclass[11pt, a4paper]{article}
\usepackage[utf8]{inputenc}
\usepackage[T1]{fontenc}
\usepackage{amsmath}
\usepackage{amsfonts}
\usepackage{amssymb}
\usepackage{geometry}
\usepackage{xcolor}
\usepackage{hyperref}
\usepackage{verbatim}
\usepackage{booktabs}

% Configurazione dei margini per prevenire overflow
\geometry{
    a4paper,
    margin=1in,
}

% Definisco un colore base per i titoli (es. blu)
\definecolor{MyBlue}{rgb}{0.1, 0.1, 0.6}

\begin{document}

\color{MyBlue}
\section*{Architetture dei Sistemi di Elaborazione \texttt{02GOLOV}}
\subsection*{Computer Architectures \texttt{02LSEYG}}
\subsubsection*{Laboratory 0x02}
\color{black}

\textbf{Expected delivery of \texttt{lab\_02.zip} including:}
\begin{itemize}
    \item \texttt{program\_1.s}
    \item \texttt{lab\_02.pdf} (fill and export this file to pdf)
\end{itemize}
\textbf{Delivery date 23/10/2025}

All the previous steps for compiling the source code, simulating with \texttt{gem5} and visualizing the pipeline have been collected into a workspace available at:
\texttt{\href{https://github.com/cad-polito-it/ase\_riscv\_gem5\_sim}{https://github.com/cad-polito-it/ase\_riscv\_gem5\_sim}}
To create your simulation workspace, you can run the following git commands:
\begin{itemize}
    \item For HTTPS clone: \texttt{\textasciitilde/my\_gem5Dir\$ git clone https://github.com/cad-polito-it/ase\_riscv\_gem5\_sim.git}
    \item For SSH: \texttt{\textasciitilde/my\_gem5Dir\$ git clone git@github.com:cad-polito-it/ase\_riscv\_gem5\_sim.git}
\end{itemize}
If you wish to install on your machine the tools (without using the VM or LABINF’s PCs) the repository contains information (and scripts) on installing the necessary tools out of the box (and generate a \texttt{setup\_default} file accordingly). See README.

\textbf{IMPORTANT:} The repository contains three \texttt{setup\_default} files:
\begin{itemize}
    \item \texttt{setup\_default}: an example of out of box installation and different tool paths.
    \item \texttt{setup\_default.vm}: tool paths used in the virtual machines.
    \item \texttt{setup\_default.labinf}: tool paths used in the LABINF machines.
\end{itemize}
Follow the HOWTO instructions available on the GitHub Repository for simulating a program. You simply run the following command: \texttt{\textasciitilde/ase\_riscv\_gem5\_sim\$ ./simulate.sh}
And select the program you want to simulate and see on the pipeline visualizer.
A useful list of RISC-V instructions can be found at:
\texttt{\href{https://msyksphinz-self.github.io/riscv-isadoc/\#\_rv32i\_rv64i\_instructions}{https://msyksphinz-self.github.io/riscv-isadoc/\#\_rv32i\_rv64i\_instructions}}

\bigskip
\color{MyBlue}
\subsection*{\#\#\#\# - Exercise 1}
\color{black}
The \texttt{riscv\_in\_order\_hen\_patt.py} (in the \texttt{gem5} folder inside the workspace) file contains the pipeline configuration. Here you can modify the operation latency and the issue latency of integer and floating-point functional units (ALU, Multiplier, Divider).
You must configure the Gem5 simulator with the Initial Configuration as described below:
\begin{verbatim}
INTEGER_ALU_LATENCY = 1  INTEGER_MUL_LATENCY = 1
INTEGER_DIV_LATENCY = 1  FLOAT_ALU_LATENCY = 8
FLOAT_MUL_LATENCY = 24  FLOAT_DIV_LATENCY = 42
\end{verbatim}

Write an assembly program (\texttt{program\_1.s}) for the RISCV architecture described before being able to implement the following high-level code:
\begin{verbatim}
for (i = 31; i >= 0; i--) {
  v4[i] = v1[i]*v1[i] – v2[i];
  v5[i] = v4[i]/v3[i] – v2[i];
  v6[i] = (v4[i]-v1[i])*v5[i];
}
\end{verbatim}
Assume that the vectors v1[], v2[], and v3[] have been previously allocated in memory and contain 32 single-precision floating-point values; also assume that v3[] does not contain any ‘0’. Additionally, the vectors v4[], v5[], v6[] are empty vectors allocated in memory. Calculate the data memory footprint of your program:

\subsubsection*{Data Memory Footprint}
\begin{table}[h]
\centering
\begin{tabular}{|l|c|}
\hline
\textbf{Data} & \textbf{Number of bytes} \\ \hline
V1 & 128 \\ \hline
V2 & 128 \\ \hline
V3 & 128 \\ \hline
V4 & 128 \\ \hline
V5 & 128 \\ \hline
V6 & 128 \\ \hline
\textbf{Total} & \textbf{768} \\ \hline
\end{tabular}
\end{table}

\textbf{Calculate the CPU performance equation (CPU time) of the above program by assuming a clock frequency of 15 MHz:}
\[
\text{CPU time} = \left(\sum_{i=1}^{n} \text{CPI}_i \times \text{IC}_i\right) \times \text{Clock cycle period}
\]
By definition:
\begin{itemize}
    \item $\text{CPI}_i$ is equal to the number of clock cycles required by the related functional unit to execute the instruction (EX stage);
    \item $\text{IC}_i$ is the number of times an instruction is repeated in the referenced source code.
\end{itemize}

\textbf{Recalculate the CPU performance equation assuming that you can triple the speed of just one unit at a time:}
\begin{itemize}
    \item b) FP multiplier (MUL) unit latency: $24 \rightarrow 8$ clock cycles
    \item c) FP divider (DIV) unit latency: $42 \rightarrow 14$ clock cycles
\end{itemize}

\begin{table}[h]
\centering
\caption{Calculate the CPU time by hand}
\label{tab:cpu_time_hand}
\begin{tabular}{|l|c|c|c|}
\hline
& \textbf{Initial CPU time (a)} & \textbf{CPU time (b – MUL speeded up)} & \textbf{CPU time (c – DIV speeded up)} \\ \hline
\texttt{program\_1.s} & 0.0002807 s & 0.0002125 s & 0.0001528 s \\ \hline
\end{tabular}
\end{table}

\textbf{Using the simulator, calculate the CPU time again and fill in the following table:}
\begin{table}[h]
\centering
\caption{Collect the CPU time using the simulator}
\label{tab:cpu_time_sim}
\begin{tabular}{|l|c|c|c|}
\hline
& \textbf{Initial CPU time (a)} & \textbf{CPU time (b – MUL speeded up)} & \textbf{CPU time (c – DIV speeded up)} \\ \hline
\texttt{program\_1.s} & 0.0002892 s & 0.000221 s & 0.0001612 s \\ \hline
\end{tabular}
\end{table}

Are there any differences? If so, where and why? If not, please provide some comments in the box below:

\textbf{Your answer:} Yes, there are notable differences between simulated CPU times and those calculated by hand. This is because theoretical formula gives us the nominal CPU time, assuming that each instruction is completely executed in a fixed number of cycles without interruptions. During the simulation, due to hazards, stalls and memory delays, the actual CPU time increases reducing the performance.

\newpage
\color{MyBlue}
\subsection*{\#\#\#\# - Exercise 2}
\color{black}
Using the Gem5 simulator, validate experimentally the Amdahl’s law, defined as follows:
\begin{itemize}
    \item a. Using the program developed before \texttt{program\_1.s}
    \item b. Modify the processor architectural parameters related (in \texttt{riscv\_in\_order\_hen\_patt.py}) to multicycle instructions latency in the following way:
\end{itemize}
\[
\text{speedup}_{\text{overall}} = \frac{\text{execution time}_{\text{old}}}{\text{execution time}_{\text{new}}} = \frac{1}{(1 - \text{fraction}) + \frac{\text{fraction}}{\text{speedup}_{\text{enhanced}}}}
\]

\textbf{Configuration 1:}
\begin{verbatim}
INTEGER_ALU_LATENCY = 1  INTEGER_MUL_LATENCY = 1
INTEGER_DIV_LATENCY = 1  FLOAT_ALU_LATENCY = 6
FLOAT_MUL_LATENCY = 20  FLOAT_DIV_LATENCY = 38
\end{verbatim}

\textbf{Configuration 2:}
\begin{verbatim}
INTEGER_ALU_LATENCY = 1  INTEGER_MUL_LATENCY = 1
INTEGER_DIV_LATENCY = 1  FLOAT_ALU_LATENCY = 4
FLOAT_MUL_LATENCY = 16  FLOAT_DIV_LATENCY = 30
\end{verbatim}

\textbf{Configuration 3:}
\begin{verbatim}
INTEGER_ALU_LATENCY = 1  INTEGER_MUL_LATENCY = 1
INTEGER_DIV_LATENCY = 1  FLOAT_ALU_LATENCY = 2
FLOAT_MUL_LATENCY = 4  FLOAT_DIV_LATENCY = 8
\end{verbatim}

Compute both manually (using the Amdahl’s Law) and with the simulator the speed-up for any one of the previous processor configurations. Compare the obtained results and complete the following table.

\begin{table}[h]
\centering
\caption{\texttt{program\_1.s} speed-up computed by hand and by simulation}
\label{tab:speedup}
\begin{tabular}{|l|c|c|c|c|}
\hline
\textbf{Proc. Config.} & \textbf{Initial config. [c.c.]} & \textbf{Config. 1 [c.c.]} & \textbf{Config. 2 [c.c.]} & \textbf{Config. 3 [c.c.]} \\ \hline
By hand & 1 & 1.584 & 1.5222 & 3.3163 \\ \hline
By simulation & 1 & 1.1534 & 1.4195 & 3.1233 \\ \hline
\end{tabular}
\end{table}

\end{document}
